\documentclass[11pt]{article}
\usepackage[margin=1in]{geometry}
\usepackage{xcolor}
\usepackage{tcolorbox}
\usepackage{hyperref}
\usepackage{enumitem}
\setlist{noitemsep,topsep=0pt}
\usepackage{booktabs}
\usepackage{listings}
\usepackage{amsmath}
\usepackage{amssymb}

% Define OSU orange and AI blue colors
\definecolor{osaorange}{RGB}{220,115,38}
\definecolor{aiblue}{RGB}{30,144,255}

% Configure hyperref colors
\hypersetup{
    colorlinks=true,
    linkcolor=blue,
    urlcolor=blue,
    citecolor=blue
}

% Configure R code listings
\lstset{
    language=R,
    basicstyle=\ttfamily\small,
    breaklines=true,
    frame=single,
    backgroundcolor=\color{gray!10},
    numbers=left,
    numberstyle=\tiny,
    keywordstyle=\color{blue},
    commentstyle=\color{green!60!black},
    stringstyle=\color{red}
}

\title{}
\author{}
\date{}

\begin{document}

% Orange header box with black outline
\begin{tcolorbox}[colback=osaorange,colframe=black,boxrule=2pt,arc=0pt,left=6pt,right=6pt,top=10pt,bottom=10pt]
\centering
\textcolor{white}{\textbf{\Large In-Class Worksheet}}\\
\textcolor{white}{\textbf{\large Data Visualization and R Programming}}
\end{tcolorbox}

\vspace{8pt}

\noindent\textbf{Name:} \rule{8cm}{0.4pt} \hfill \textbf{Date:} \rule{4cm}{0.4pt}\\[0.5ex]
\textbf{Course:} H524 Introduction to Biostatistics\\
\textbf{Topic:} Data Presentation, Histograms, Boxplots, Introduction to R Programming

\vspace{8pt}

\section{Part 1: Review and Setup}

\subsection{Exercise 1.1: Quick Review}
From previous material, what were the three measures of central tendency for our blood pressure data \texttt{c(170, 193, 110, 135, 135, 115)}?

\begin{itemize}
\item Mean: \rule{2cm}{0.4pt} mmHg
\item Median: \rule{2cm}{0.4pt} mmHg
\item Mode: \rule{2cm}{0.4pt} mmHg
\end{itemize}

\begin{lstlisting}
# Load the data again
blood_pressure <- c(170, 193, 110, 135, 135, 115)

# Verify your answers
mean(blood_pressure)
median(blood_pressure)
\end{lstlisting}

\section{Part 2: Introduction to Data Visualization}

\subsection{Exercise 2.1: Creating Histograms}

We'll work with simulated guinea pig survival data from three groups:

\begin{lstlisting}
# Guinea pig survival data (tuberculosis study)
control <- c(312, 345, 287, 398, 421, 356, 389, 334, 367, 401)
low_dose <- c(289, 267, 301, 234, 198, 245, 278, 213, 256, 298)
high_dose <- c(156, 189, 134, 167, 123, 145, 178, 201, 143, 134)

# Create histograms
hist(control, main="Control Group Survival", xlab="Days", col="lightblue")
hist(low_dose, main="Low Dose Group Survival", xlab="Days", col="lightgreen")
hist(high_dose, main="High Dose Group Survival", xlab="Days", col="lightcoral")
\end{lstlisting}

\textbf{Key Questions:}
\begin{enumerate}
\item What do you notice about the shape of each distribution?
\item Which group appears to have the best survival prospects?
\item Why is it important to use the same x-axis scale when comparing groups?
\end{enumerate}

\subsection{Exercise 2.2: AI-Enhanced Histogram Creation}

\begin{tcolorbox}[colback=aiblue!10,colframe=aiblue,boxrule=1pt,arc=2pt,left=6pt,right=6pt,top=6pt,bottom=6pt]
\textbf{\textcolor{aiblue}{AI Prompt:}} ``Help me create side-by-side histograms for comparing three groups of guinea pig survival data. The data should use the same scales and bin widths for fair comparison.''
\end{tcolorbox}

\textbf{AI Response (fill in):}
\begin{lstlisting}
# Paste AI-generated code here:


\end{lstlisting}

\textbf{Verification Checklist:}
\begin{itemize}
\item[$\square$] Same x-axis scale across all histograms?
\item[$\square$] Same number of bins?
\item[$\square$] Clear labels and titles?
\item[$\square$] Different colors for each group?
\end{itemize}

\section{Part 3: Boxplots for Group Comparison}

\subsection{Exercise 3.1: Understanding the Five-Number Summary}

For the control group data, calculate the five-number summary:

\begin{lstlisting}
control <- c(312, 345, 287, 398, 421, 356, 389, 334, 367, 401)

# Calculate five-number summary
summary(control)

# Individual components
min(control)
quantile(control, 0.25)  # Q1
median(control)          # Q2
quantile(control, 0.75)  # Q3
max(control)
\end{lstlisting}

\textbf{Fill in the five-number summary:}
\begin{itemize}
\item Minimum: \rule{2cm}{0.4pt}
\item Q1 (25th percentile): \rule{2cm}{0.4pt}
\item Q2 (Median): \rule{2cm}{0.4pt}
\item Q3 (75th percentile): \rule{2cm}{0.4pt}
\item Maximum: \rule{2cm}{0.4pt}
\end{itemize}

\subsection{Exercise 3.2: Creating Boxplots}

\begin{lstlisting}
# Create individual boxplots
boxplot(control, main="Control Group", ylab="Survival Days")
boxplot(low_dose, main="Low Dose Group", ylab="Survival Days")
boxplot(high_dose, main="High Dose Group", ylab="Survival Days")

# Side-by-side comparison
boxplot(control, low_dose, high_dose,
        names=c("Control", "Low Dose", "High Dose"),
        main="Guinea Pig Survival by Treatment Group",
        ylab="Survival Days",
        col=c("lightblue", "lightgreen", "lightcoral"))
\end{lstlisting}

\textbf{Analysis Questions:}
\begin{enumerate}
\item Which group has the highest median survival?
\item Which group shows the most variability (largest IQR)?
\item Are there any potential outliers visible?
\end{enumerate}

\section{Part 4: R Programming Fundamentals}

\subsection{Exercise 4.1: Working with Vectors and Data Frames}

\begin{lstlisting}
# Create a data frame combining all groups
survival_data <- data.frame(
  days = c(control, low_dose, high_dose),
  group = c(rep("Control", length(control)),
            rep("Low Dose", length(low_dose)),
            rep("High Dose", length(high_dose)))
)

# Explore the data frame
str(survival_data)
head(survival_data)
summary(survival_data)
\end{lstlisting}

\subsection{Exercise 4.2: Basic Data Manipulation}

\begin{lstlisting}
# Calculate group means
aggregate(days ~ group, data = survival_data, FUN = mean)

# Calculate group standard deviations
aggregate(days ~ group, data = survival_data, FUN = sd)

# Create a subset for just the control group
control_subset <- subset(survival_data, group == "Control")
\end{lstlisting}

\textbf{Fill in the results:}
\begin{itemize}
\item Control group mean: \rule{3cm}{0.4pt}
\item Low dose group mean: \rule{3cm}{0.4pt}
\item High dose group mean: \rule{3cm}{0.4pt}
\end{itemize}

\section{Part 5: AI-Enhanced Analysis Challenge}

\subsection{Exercise 5.1: Complete Statistical Summary}

\begin{tcolorbox}[colback=aiblue!10,colframe=aiblue,boxrule=1pt,arc=2pt,left=6pt,right=6pt,top=6pt,bottom=6pt]
\textbf{\textcolor{aiblue}{AI Challenge:}} Use AI to help create a comprehensive analysis comparing the three treatment groups.

\textbf{Your AI Prompt:}
``I have guinea pig survival data for three groups: control, low dose, and high dose. Help me create a complete statistical comparison including descriptive statistics, appropriate visualizations, and interpretation of the results.''
\end{tcolorbox}

\textbf{AI Response Documentation:}
\begin{enumerate}
\item \textbf{Descriptive Statistics Code:}
\begin{lstlisting}
# Paste AI code here:


\end{lstlisting}

\item \textbf{Visualization Code:}
\begin{lstlisting}
# Paste AI code here:


\end{lstlisting}

\item \textbf{AI Interpretation:}
(Write the AI's interpretation of the results)

\vspace{3cm}

\end{enumerate}

\textbf{Your Verification:}
\begin{itemize}
\item[$\square$] Are the statistical calculations correct?
\item[$\square$] Do the visualizations effectively show group differences?
\item[$\square$] Is the interpretation scientifically sound?
\item[$\square$] What would you add or modify?
\end{itemize}

\section{Part 6: Connecting to Biostatistics Concepts}

\subsection{Exercise 6.1: Population vs. Sample Discussion}

\textbf{Questions for Discussion:}
\begin{enumerate}
\item In this guinea pig study, what is the population of interest?
\item What is our sample?
\item How might these results generalize to the broader population?
\item What are potential limitations of this study design?
\end{enumerate}

\subsection{Exercise 6.2: Setting Up for Statistical Inference}

In upcoming classes we'll learn about probability. Based on today's data:

\begin{enumerate}
\item If we randomly selected one guinea pig from the control group, would you expect its survival to be exactly the mean survival time? Why or why not?

\item How confident are you that the control group truly has better survival than the high dose group? What additional information would help you be more certain?
\end{enumerate}

\section{Wrap-up and Preview}

\textbf{Key Takeaways:}
\begin{itemize}
\item Histograms show the shape and distribution of data
\item Boxplots efficiently compare groups using five-number summaries
\item R data frames organize different types of data together
\item AI can accelerate analysis but requires human verification
\end{itemize}

\textbf{Next Class:} We'll dive deeper into numerical summary measures, rates and standardization, and advanced R programming techniques.

\section{AI Usage Documentation}
\begin{tcolorbox}[colback=gray!10,colframe=gray,boxrule=1pt,arc=2pt,left=6pt,right=6pt,top=6pt,bottom=6pt]
\textbf{Document your AI interaction:}
\begin{itemize}
\item AI tools used: \rule{8cm}{0.4pt}
\item Most helpful AI suggestion: \rule{8cm}{0.4pt}
\item Code that needed correction: \rule{8cm}{0.4pt}
\item Verification methods used: \rule{8cm}{0.4pt}
\item Overall AI effectiveness (1-10): \rule{2cm}{0.4pt}
\end{itemize}
\end{tcolorbox}

\end{document}